% Part: first-order-logic
% Chapter: completeness
% Section: compactness

\documentclass[../../../include/open-logic-section]{subfiles}

\begin{document}

\iftag{FOL}
      {\olfileid{fol}{com}{com}}
      {\olfileid{pl}{com}{com}}

\olsection{సంహతత్వ సిద్ధాంతం}

సంపూర్ణతా సిద్ధాంతపు ఒక ముఖ్యమైన పరిణామం సంహతత్వ సిద్ధాంతం.
!!{sentence}s సమితిలోని ప్రతి
\emph{పరిమిత} ఉపసమితి సంతృప్తిపరచదగినదైతే, మొత్తం సమితి---అది
అనంతమైనదైనా సరే---సంతృప్తిపరచదగినదని సంహతత్వ సిద్ధాంతం చెబుతుంది.
ఇది ఏమాత్రం స్పష్టమైన విషయం కాదు. మొదటి చూపులోనైనా, వైరుధ్యం అనంతమైన
సంఖ్య వల్లనే ఉత్పన్నమయ్యే వైరుధ్యభరిత అనంత !!{sentence}s సమితులు
ఉండే అవకాశాన్ని ఏదీ తొలగిస్తున్నట్లు కనిపించదు. అలాంటి పరిస్థితిని
తొలగించవచ్చని సంహతత్వ సిద్ధాంతం చెబుతుంది: ప్రతి పరిమిత ఉపసమితీ
సంతృప్తిపరచదగినదై ఉండే, కానీ మొత్తం మీద సంతృప్తిపరచలేని అనంత
!!{sentence}s సమితులు లేవు. సంపూర్ణతా సిద్ధాంతంలాగే, దీనికి అనుగమనానికి
సంబంధించిన ఒక రూపం ఉంది: ఒక అనంత !!{sentence}s సమితి ఏదైనా విషయాన్ని
అనుగమింపజేస్తే, ఇప్పటికే దాని ఒక పరిమిత ఉపసమితి దాన్ని అనుగమింపజేస్తుంది.

\begin{defn}
  !!{formula}s సమితి~$\Gamma$లోని ప్రతి పరిమిత $\Gamma_0 \subseteq
  \Gamma$ సంతృప్తిపరచదగినదైతే, అప్పుడు మాత్రమే ఆ సమితిని
  \emph{పరిమితంగా సంతృప్తిపరచదగినది} అంటాం.
\end{defn}

\begin{thm}[సంహతత్వ సిద్ధాంతం]
\ollabel{thm:compactness}
ఏ వాక్యాల సమితి~$\Gamma$కు, ఏ వాక్యం~$!A$కు కిందివి
వర్తిస్తాయి:
\sourcecorrection{OLTECOM-008}{ఆంగ్ల మూలం $\Gamma$, $!A$ రెండింటినీ
``వాక్యాలు'' అని ఒకే రకంగా పేర్కొంది. తరువాతి రెండు అంశాలు కోరినట్లుగా
$\Gamma$ను వాక్యాల సమితిగా, $!A$ను ఒక వాక్యంగా స్పష్టంచేశాం.}
\begin{enumerate}
  \item $\Gamma \Entails !A$ అయినప్పుడు, అప్పుడు మాత్రమే $\Gamma_0
    \subseteq \Gamma$, $\Gamma_0 \Entails !A$ అయ్యే పరిమిత ఉపసమితి
    ఉంది.
  \item $\Gamma$ సంతృప్తిపరచదగినదైతే, అప్పుడు మాత్రమే అది పరిమితంగా
    సంతృప్తిపరచదగినది.
\end{enumerate}
\end{thm}

\begin{proof}
(2)ను నిరూపిస్తాం. $\Gamma$ సంతృప్తిపరచదగినదైతే, అన్ని $!A \in
\Gamma$లకు $\iftag{FOL}{\Sat{M}}{\pSat{v}}{!A}$ అయ్యే
\iftag{FOL}{!!a{structure}~$\Struct{M}$}{!!a{valuation}~$\pAssign{v}$}
ఉంటుంది. ఈ $\iftag{FOL}{\Struct{M}}{\pAssign{v}}$ సహజంగానే
~$\Gamma$లోని ప్రతి పరిమిత ఉపసమితిని కూడా సంతృప్తిపరుస్తుంది; కాబట్టి
$\Gamma$ పరిమితంగా సంతృప్తిపరచదగినది.

ఇప్పుడు~$\Gamma$ పరిమితంగా సంతృప్తిపరచదగినదని అనుకుందాం. అప్పుడు ప్రతి
పరిమిత ఉపసమితి~$\Gamma_0 \subseteq \Gamma$ సంతృప్తిపరచదగినది.
నిర్దుష్టత ప్రకారం
(\iftag{FOL}{%
  \tagrefs{prfAX/{fol:axd:sou:cor:consistency-soundness},
    prfSC/{fol:seq:sou:cor:consistency-soundness},
    prfND/{fol:ntd:sou:cor:consistency-soundness},
    prfTab/{fol:tab:sou:cor:consistency-soundness}}}{%
  \tagrefs{prfAX/{pl:axd:sou:cor:consistency-soundness},
    prfSC/{pl:seq:sou:cor:consistency-soundness},
    prfND/{pl:ntd:sou:cor:consistency-soundness},
    prfTab/{pl:tab:sou:cor:consistency-soundness}}}),
ప్రతి పరిమిత ఉపసమితీ అవైరుధ్యమైనది. అప్పుడు కింది ఫలితం ప్రకారం
$\Gamma$ స్వయంగా అవైరుధ్యమైనది:
\iftag{FOL}{%
  \tagrefs{prfAX/{fol:axd:ptn:prop:proves-compact},
    prfSC/{fol:seq:ptn:prop:proves-compact},
    prfND/{fol:ntd:ptn:prop:proves-compact},
    prfTab/{fol:tab:ptn:prop:proves-compact}}}{%
  \tagrefs{prfAX/{pl:axd:ptn:prop:proves-compact},
    prfSC/{pl:seq:ptn:prop:proves-compact},
    prfND/{pl:ntd:ptn:prop:proves-compact},
    prfTab/{pl:tab:ptn:prop:proves-compact}}}.
సంపూర్ణత ప్రకారం (\olref[cth]{thm:completeness}), $\Gamma$
అవైరుధ్యమైనది కనుక అది సంతృప్తిపరచదగినది.
\end{proof}

\tagprob{FOL}
\begin{prob}
\olref[fol][com][com]{thm:compactness}లోని (1)ను నిరూపించండి.
\end{prob}
\tagendprob

\tagprob{notFOL}
\begin{prob}
\olref[pl][com][com]{thm:compactness}లోని (1)ను నిరూపించండి.
\end{prob}
\tagendprob

\iftag{FOL}{%
\begin{ex}
సిద్ధాంతం~$\Gamma$ యొక్క ప్రతి నమూనా~$\Struct{M}$లో, ప్రతి పదం~$t$
సహజంగానే~$\Domain{M}$లోని !!a{element}ను సూచిస్తుంది. $\Domain{M}$లోని
ప్రతి !!{element}ను కూడా ఏదో ఒక పదం సూచిస్తుందని హామీ ఇవ్వగలమా?
మరోలా చెప్పాలంటే, అన్ని నమూనాలూ ఆవరితంగా ఉండే సిద్ధాంతాలు~$\Gamma$
ఉన్నాయా? $\Gamma$కు అనంత నమూనాలు ఉంటే అలా కాదని సంహతత్వ సిద్ధాంతం
చూపుతుంది. అది ఇలా: $\Struct{M}$ అనేది~$\Gamma$ యొక్క అనంత నమూనా,
$c$ అనేది~$\Gamma$ భాషలో లేని !!a{constant} అనుకుందాం. $t$ అనేది
~$\Gamma$ భాష~$\Lang{L}$లోని పదంగా ఉండే అన్ని
$\eq/[c][t]$ వాక్యాల సమితిని~$\Delta$ అందాం; అంటే,
  \[
  \Delta = \Setabs{\eq/[c][t]}{t \in \Trm[L]}.
  \]
$\Gamma \cup \Delta$లోని ఒక పరిమిత ఉపసమితిని $\Gamma' \cup
\Delta'$గా రాయవచ్చు; ఇక్కడ $\Gamma' \subseteq \Gamma$, $\Delta'
\subseteq \Delta$. $\Delta'$ పరిమితమైనది కనుక, దానిలో పరిమిత
సంఖ్యలో పదాలు మాత్రమే ఉండగలవు. వాటిలో ఏదీ సూచించని
$a \in \Domain{M}$ అనే ఒక $\Domain{M}$లోని !!a{element}ను తీసుకొని,
$\Struct{M}$లాగే
ఉండి అదనంగా $\Assign{c}{M'} = a$ అయ్యే !!{structure}ను~$\Struct{M'}$
అందాం. $\Delta'$లో సంభవించే అన్ని~$t$లకు $a \neq \Value{t}{M}$ కనుక
$\Sat{M'}{\Delta'}$. $\Sat{M}{\Gamma}$, $\Gamma' \subseteq \Gamma$,
$c$ అనేది~$\Gamma$లో సంభవించదు కనుక $\Sat{M'}{\Gamma'}$ కూడా.
కలిపితే, $\Gamma \cup \Delta$లోని ప్రతి పరిమిత ఉపసమితి $\Gamma'
\cup \Delta'$కు $\Sat{M'}{\Gamma' \cup \Delta'}$. కాబట్టి $\Gamma
\cup \Delta$లోని ప్రతి పరిమిత ఉపసమితి సంతృప్తిపరచదగినది. సంహతత్వం
ప్రకారం $\Gamma \cup \Delta$ స్వయంగా సంతృప్తిపరచదగినది. అందువల్ల
$\Sat{M}{\Gamma \cup \Delta}$ అయ్యే నమూనాలు ఉన్నాయి. అలాంటి ప్రతి
$\Struct{M}$, $\Gamma$కు ఒక నమూనా; కానీ~$\Lang{L}$లోని అన్ని పదాలు~$t$కు
$\Value{c}{M} \neq \Value{t}{M}$ కనుక అది ఆవరితమైనది కాదు.
\end{ex}

\begin{ex}
!!{predicate}~$<$, !!{constant}s~$\Obj{0}$, $\Obj{1}$, ఇంకా
!!{function}s~$+$, $\times$, $-$లను కలిగిన !!a{language}~$\Lang{L}$ను
పరిశీలించండి. ప్రదేశం~$\Rat$, స్పష్టమైన అర్థవివరణలు గల
!!{structure}~$\Struct{Q}$లో సత్యమైన ఈ !!{language}లోని అన్ని
!!{sentence}s సమితిని~$\Gamma$ అందాం. కరణీయ సంఖ్యల గురించి సత్యమైన
~$\Lang{L}$లోని అన్ని !!{sentence}s సమితి~$\Gamma$. సహజంగానే
~$\Rat$లో (అంతేకాదు~$\Real$లో కూడా), $0$ కంటే పెద్దదై, ప్రతి
$k \in \PosInt$కు $1/k$ కంటే చిన్నదైన సంఖ్య~$r$ ఏదీ లేదు. అలాంటి
సంఖ్య ఉంటే అది
\emph{అనంతసూక్ష్మం:} సున్నా కాదు, కానీ అనంతంగా చిన్నది. అనంతసూక్ష్మాలు
ఉండే~$\Gamma$ నమూనాలు ఉన్నాయని చూపడానికి సంహతత్వ సిద్ధాంతాన్ని
ఉపయోగించవచ్చు. మన భాషలో భాగహారానికి !!a{function} లేదు (సున్నాతో
భాగహారం నిర్వచితం కాదు; !!{function}sను సర్వత్ర నిర్వచిత ప్రమేయాలతోనే
అర్థం చేసుకోవాలి). అయినా $r < 1/k$ అని వ్యక్తం చేయవచ్చు; ఎందుకంటే
$r \cdot k < 1$ అయినప్పుడు, అప్పుడు మాత్రమే అది నిజం. ఇప్పుడు $c$ ఒక
కొత్త !!{constant} అనుకొని, $\Delta$ను ఇలా ఉంచండి:
\[
\{\Obj{0} < c\}
\cup \Setabs{ c \times \num k < \Obj{1} }{k \in \PosInt}
\]
(ఇక్కడ $\num{k} = (\Obj{1} + (\Obj{1} + \dots + (\Obj{1} +
\Obj{1})\dots))$; అందులో $k$~$\Obj{1}$లు ఉంటాయి). $\Delta$లోని ఏ
పరిమిత ఉపసమితి~$\Delta_0$కైనా, $\Delta_0$లో ఉన్న అన్ని
!!{sentence}s $c \times \num{k} < \Obj{1}$కు $k < K$ అయ్యే ఒక~$K$
ఉంటుంది. $\Assign{c}{Q'} = 1/K$గా ఉంచి $\Struct{Q}$ను~$\Struct{Q'}$గా
విస్తరిస్తే, ఏ పరిమిత $\Gamma_0 \subseteq \Gamma$కైనా
$\Sat{Q'}{\Gamma_0 \cup \Delta_0}$. అందువల్ల $\Gamma \cup \Delta$
పరిమితంగా సంతృప్తిపరచదగినది (దీన్ని వివరంగా నిరూపించడం అభ్యాసం).
సంహతత్వం ప్రకారం $\Gamma \cup \Delta$ సంతృప్తిపరచదగినది. $\Gamma
\cup \Delta$ యొక్క ఏ నమూనా~$\Struct{S}$లోనైనా ఒక అనంతసూక్ష్మం ఉంది;
అది~$\Assign{c}{S}$.
\end{ex}
}{}

\tagprob{FOL}
\begin{prob}
అంకగణితపు ప్రామాణిక నమూనా~$\Struct{N}$లో, ప్రతి సూత్రం~$\num{n} < x$
ను సంతృప్తిపరిచే !!{element}~$k \in \Domain{N}$ ఏదీ లేదు (ఇక్కడ
$\num{n}$ అనేది $n$~$\prime$లతో $\Obj{0}^{\prime\dots\prime}$).
అంకగణితపు ప్రామాణిక నమూనా~$\Struct{N}$లో సత్యమైన అంకగణిత భాషలోని
!!{sentence}s, ప్రతి సూత్రం~$\num{n} < x$ను
\emph{సంతృప్తిపరిచే} !!a{element} ఉన్న !!a{structure}~$\Struct{N'}$లో
కూడా సత్యమని చూపడానికి సంహతత్వ సిద్ధాంతాన్ని ఉపయోగించండి.
\end{prob}
\tagendprob

\iftag{FOL}{
\begin{ex}
!!{identity} గల ప్రథమక్రమ తర్కం, !!{domain} పరిమాణానికి ఏదో ఒక కనిష్ఠ
పరిమాణం తప్పనిసరి అని వ్యక్తం చేయగలదని మనకు తెలుసు: $!A_{\ge n}$ అనే
వాక్యం (``కనీసం $n$ వేర్వేరు వస్తువులు ఉన్నాయి'' అని అది చెబుతుంది)
$\Domain{M}$లో కనీసం~$n$ వస్తువులు గల నిర్మాణాల్లో మాత్రమే సత్యం.
కాబట్టి
  \[
  \Delta = \Setabs{!A_{\ge n}}{n \ge 1}
  \]
గా తీసుకుంటే, $\Delta$ యొక్క ఏ నమూనా అయినా అనంతంగా ఉండాలి. అందువల్ల
~$\Delta$ను ఒక సిద్ధాంతానికి చేర్చడం ద్వారా దానికి అనంత నమూనాలు మాత్రమే
ఉండేలా హామీ ఇవ్వవచ్చు: $\Gamma \cup \Delta$ నమూనాలు సరిగ్గా~$\Gamma$
యొక్క అనంత నమూనాలే.

కాబట్టి ప్రథమక్రమ తర్కం అనంతత్వాన్ని వ్యక్తం చేయగలదు. అయితే అది
పరిమితత్వాన్ని వ్యక్తం చేయలేదని సంహతత్వ సిద్ధాంతం చూపుతుంది. ఏదో ఒక
వాక్యాల సమితి~$\Lambda$ను అన్ని పరిమిత !!{structure}s, అవి మాత్రమే,
సంతృప్తిపరుస్తాయని అనుకుందాం. అప్పుడు $\Delta \cup \Lambda$ పరిమితంగా
సంతృప్తిపరచదగినది. ఎందుకు? $\Delta' \subseteq \Delta$, $\Lambda'
\subseteq \Lambda$ అయ్యే పరిమిత $\Delta' \cup \Lambda' \subseteq
\Delta \cup \Lambda$ను తీసుకోండి. $!A_{\ge n} \in \Delta'$ అయ్యే
అతిపెద్ద సంఖ్యను~$n$ అందాం. అన్ని పరిమిత నిర్మాణాలు~$\Lambda$ను
సంతృప్తిపరుస్తాయి కనుక, పరిమిత సంఖ్యలో, కానీ $\ge n$ !!{element}s గల
నమూనా~$\Struct{M}$ దానికి ఉంది. అప్పుడు $\Sat{M}{\Delta' \cup
\Lambda'}$. సంహతత్వం ప్రకారం $\Delta \cup \Lambda$కు ఒక అనంత నమూనా
ఉంది; ఇది~$\Lambda$ను పరిమిత !!{structure}s మాత్రమే సంతృప్తిపరుస్తాయనే
పరికల్పనకు విరుద్ధం.
\end{ex}
}{}

\end{document}
